A mudança exata de área é
\[
(1+\alpha\Delta T)^{2}-1 = 2\alpha\Delta T + (\alpha\Delta T)^{2}.
\]
O termo negligenciado ao usar apenas o coeficiente de superfície é $(\alpha\Delta T)^{2}$. Para que esse termo seja 1% da mudança de área, basta escolher $\alpha\Delta T \approx 0{,}02$. Com $\alpha\approx2{,}4\times10^{-5}\ \mathrm{K}^{-1}$ para o alumínio,
\[
\Delta T \approx \frac{0{,}02}{2{,}4\times10^{-5}} \approx 8,3\times10^{2}\ \mathrm{K}.
\]